% Tóm tắt (tiếng Việt, một trang): nội dung tương ứng Abstract.
\FrontChapter{TÓM TẮT}
\begin{otherlanguage}{vietnamese}
Một phép thay thế amino acid có thể làm thay đổi hoạt tính degron, tín hiệu trình tự liên quan đến
phân hủy protein có chọn lọc. Luận văn dự đoán hiệu ứng đột biến bằng độ thay đổi chỉ số ổn định
protein ($\Delta$PSI) trong thí nghiệm reporter peptide và xây dựng công cụ phân tích local.
Quy trình dựng lại nhãn, kiểm soát chồng lặp trình tự và đánh giá theo nhóm protein, với lựa chọn
mô hình ở tập trong và khoảng tin cậy bootstrap theo nhóm.

Bộ phát triển toàn dải PSI có \nVthreeNPeptides{} peptide, \nVthreeNCells{} ô đột biến và
\nVthreeNGroups{} nhóm protein. Mở rộng train ridge từ PSI hoang dã không quá 4 sang toàn dải
giảm MAE trung bình theo peptide từ \nVthreeRidgeB{} xuống \nVthreeRidgeD{} PSI.
ESM2-35M đóng băng kết hợp CNN đạt MAE \nVthreeCnnD{}, so với \nVthreePointwiseD{} của
đầu pointwise đối chứng. Tiền huấn luyện bằng nhãn scanning cải thiện CNN one-hot trong thiết kế
đã thử; trọng số chất lượng làm ridge kém hơn. Thứ hạng MAE chưa xác lập kiến trúc thắng.
Các thí nghiệm miền hạn chế trước đó cung cấp baseline ridge đã hiệu chuẩn và kiểm khớp giữa
đường suy luận với runner nghiên cứu.

Đến ngày 03/11/2026, đã có baseline React/FastAPI; lưu dự án, job peptide/protein và đối chiếu dữ
liệu đo còn dở. Workspace mới là prototype duyệt giao diện dùng dữ liệu tổng hợp. Mở rộng pLM và
LoRA local, hiệu chuẩn toàn dải và đánh giá nguồn giữ lại chưa xong; mô hình phục vụ chưa đổi.
Các kết quả dự đoán là bằng chứng phát triển trong assay, chưa xác nhận độc lập ngoài nguồn,
hiệu ứng trên protein nguyên vẹn, nhận diện E3 cụ thể hoặc ý nghĩa lâm sàng.

\par\vspace{6bp}\noindent\textbf{Từ khóa:} degron, độ bền protein, dự đoán hiệu ứng đột biến, hệ ubiquitin--proteasome, học máy.
\end{otherlanguage}
